\chapter{Project recovery: projected curvature на допустимой \texorpdfstring{\(\bar4\to2_R\to4\)}{4bar-to-2R-to-4} chain}

Two-form/A2 trilemma закрыл три прямых маршрута, но оставил even singlet
curvature. Повторный поиск по старым superconnection, relative-projector и
vectorlike gates обнаруживает, что нужные детали уже были в проекте по
отдельности. Их объединение даёт exact determinant selector без
six-dimensional algebra node.

\section{Синтез ранее полученных механизмов}

В предыдущих разделах были независимо установлены три свойства:
\begin{enumerate}
 \item разные form degrees можно развести по ортогональным target modules;
 \item endpoint inclusion--exclusion может канонически выделять относительный
 curvature block;
 \item vectorlike fundamental modules допустимы без нового color algebra
 summand.
\end{enumerate}
Если собрать их, запрещённая шестёрка больше не нужна как node. Берутся
только обычные color \(4,\bar4\) и weak doublet \(2_R\). Шестёрка возникает
как antisymmetric двухшаговый curvature между крайними nodes.

\section{Допустимая трёхузловая chain}

Рассмотрим graded particle carrier
\begin{equation}
 E_0^+=\bar4,
 \qquad
 E_1^-=2_R,
 \qquad
 E_2^+=4.
 \label{eq:ps-valid-module-chain}
\end{equation}
Все три nodes являются fundamental или opposite modules исходной
Pati--Salam algebra. Для \(\Delta\in(2_R,4)\) зададим
\begin{equation}
 A_\Delta=\Delta:E_0\longrightarrow E_1,
 \qquad
 B_\Delta=c\,\Delta^T\varepsilon_2:E_1\longrightarrow E_2,
 \label{eq:ps-valid-module-edges}
\end{equation}
где
\begin{equation}
 \varepsilon_2=
 \begin{pmatrix}0&1\\-1&0\end{pmatrix}.
\end{equation}
SU(2) pseudoreality identity
\begin{equation}
 U_R^T\varepsilon_2=\varepsilon_2U_R^{-1}
\end{equation}
обеспечивает gauge covariance второго edge. На трёхстах random
\(SU(2)_R\times SU(4)\) transformations covariance error меньше
\(8\times10^{-15}\).

Двухшаговый path равен
\begin{equation}
 B_\Delta A_\Delta
 =c\,\Delta^T\varepsilon_2\Delta.
 \label{eq:ps-valid-module-wedge}
\end{equation}
Это antisymmetric \(4\times4\) matrix, то есть color \(6\), и
\begin{equation}
 \|B_\Delta A_\Delta\|_F^2
 =2c^2\det(\Delta\Delta^\dagger).
 \label{eq:ps-valid-module-minor-norm}
\end{equation}

\section{Почему raw trace снова недостаточен}

Self-adjoint odd operator на \(4+2+4\) carrier имеет вид
\begin{equation}
 \mathcal D_\Delta=
 \begin{pmatrix}
 0&A_\Delta^\dagger&0\\
 A_\Delta&0&B_\Delta^\dagger\\
 0&B_\Delta&0
 \end{pmatrix}.
\end{equation}
Для
\begin{equation}
 \rho=\Tr(\Delta\Delta^\dagger),
 \qquad
 \tau=\Tr[(\Delta\Delta^\dagger)^2],
 \qquad
 d=\det(\Delta\Delta^\dagger)
\end{equation}
получаем
\begin{align}
 \Tr\mathcal D_\Delta^2
 &=2(1+c^2)\rho,\\
 \Tr\mathcal D_\Delta^4
 &=2(1+c^4)\tau+8c^2d.
 \label{eq:ps-valid-module-raw-traces}
\end{align}
При каноническом \(c=1\)
\begin{equation}
 \Tr\mathcal D_\Delta^4=4\rho^2,
\end{equation}
то есть полный ordinary trace снова rank-blind. Это объясняет, почему
предыдущие minimal spectral potentials не увидели determinant, хотя нужный
path уже присутствовал внутри \(\mathcal D_\Delta^2\).

\section{Endpoint curvature projector}

Пусть \(P_0,P_2\) --- projectors на крайние irreducible nodes. Выделим
relative endpoint curvature
\begin{equation}
 F_{02}
 =P_0\mathcal D_\Delta^2P_2
 +P_2\mathcal D_\Delta^2P_0.
 \label{eq:ps-endpoint-curvature}
\end{equation}
Его два off-diagonal blocks равны
\((B_\Delta A_\Delta)^\dagger\) и \(B_\Delta A_\Delta\). Поэтому
\begin{equation}
 \boxed{
 \|F_{02}\|_F^2
 =4c^2\det(\Delta\Delta^\dagger).}
 \label{eq:ps-projected-curvature-selector}
\end{equation}
После KO6 doubling physical half-norm сохраняет ту же формулу. Machine audit
даёт maximum error меньше \(5\times10^{-13}\) при \(c=1\), а projected
curvature точно even, real-compatible и self-adjoint.

При canonical trace metrics \(c=1\) получается
\begin{equation}
 \boxed{\frac12\|F_{02}^{\rm KO6}\|_F^2
 =4\det(\Delta\Delta^\dagger).}
 \label{eq:ps-project-recovered-kappa-four}
\end{equation}
Это ровно coefficient, найденный independently path identity и combined
Casimir gap.

\section{Robustness normalization}

Если relative edge normalization пока оставить равной \(c\), induced
potential coefficient равен
\begin{equation}
 \kappa=4c^2.
\end{equation}
Предыдущий exact Hessian требует лишь
\begin{equation}
 \kappa>2
 \quad\Longleftrightarrow\quad
 c^2>\frac12.
 \label{eq:ps-project-recovery-window}
\end{equation}
Следовательно, rank-one stabilization не зависит от единственной идеально
подогнанной normalization point. Но exact numerical coefficient четыре
остаётся parent theorem только после вывода canonical equality metrics двух
edges.

\section{Что именно восстановлено из проекта}

Механизм собирает три старых результата:
\begin{enumerate}
 \item \path{s2t_v4_three_node_superconnection_closure_gate.py}: минимальный
 \((+,-,+)\) carrier и ортогональные target grades;
 \item \path{version4_relative_krajewski_star_gate.tex}: endpoint
 inclusion--exclusion и relative projectors без continuous weights;
 \item \path{version4_vectorlike_messenger_chain_gate.tex}: допустимость
 anomaly-safe fundamental-module extensions.
\end{enumerate}

\begin{proposition}[Project-recovered selector]
На допустимой fundamental-module chain \(\bar4\to2_R\to4\) relative endpoint
curvature norm автоматически равен
\(4c^2\det(\Delta\Delta^\dagger)\). При canonical module metric \(c=1\)
получается exact stabilizing coefficient четыре, тогда как полный ordinary
spectral trace остаётся rank-blind.
\end{proposition}

\section{Почему это ещё условный успех}

Найдена не ручная polynomial correction, а конкретный parent carrier и
curvature block. Тем не менее остаются четыре обязательных gates:
\begin{enumerate}
 \item вывести projector \(F\mapsto F_{02}\) из единого superconnection
 action, а не выбрать его после вычисления;
 \item вывести equality edge metrics \(c=1\) из reality, normalized trace или
 module isometry;
 \item решить, являются ли дополнительные \(4/\bar4\) physical vectorlike
 states или только auxiliary bosonic grades;
 \item пересчитать полный \((\phi,\Sigma_4,\Delta)\) Hessian и gauge-running
 ledger расширения.
\end{enumerate}

\begin{theorem}[Условное reopening]
Ранее полученные superconnection и relative-projector механизмы дают
admissible route к exact determinant selector без
six-dimensional algebra node. Но физическое замыкание требует parent-action
derivation endpoint curvature norm и normalization, поэтому status остаётся
conditional, а не closed.
\end{theorem}

\section{Последующий junk/mapping-cone audit}

Следующий gate уточняет это conditional statement. Equality edge metrics
действительно выводится: map
\(\Delta\mapsto\Delta^T\varepsilon_2\) является Frobenius-isometry, а
endpoint reflection/reality фиксирует \(c=1\) с точностью до removable
phase. Standard degree-two junk quotient, однако, не даёт projector: в
минимальном \(\mathbb C^3\) path control endpoint units \(E_{02},E_{20}\)
сами принадлежат junk. Projector выводится вместо этого как normalized
commutator с уникальной reflection-odd graph coordinate
\(h=(-1,0,1)\). Тем самым новый open gate --- не normalization и не algebraic
projection, а выбор relative mapping-cone curvature norm как физического
parent action.

\section{Литература и воспроизводимость}

Superconnection bosonic-action precedent сверён с
H.~Figueroa, J.~M.~Gracia-Bond\'ia, F.~Lizzi and J.~C.~V\'arilly,
arXiv:hep-th/9701179. Ordinary spectral action использует ungraded trace,
как в A.~H.~Chamseddine and A.~Connes, arXiv:hep-th/9606001; именно поэтому
projected superconnection curvature является дополнительной структурой, а
не автоматическим следствием raw spectral action.

Machine audit сохранён в
\path{s2t_v4_pati_salam_projected_curvature_selector.py} и
\path{s2t_v4_pati_salam_projected_curvature_selector_results.json}.